\documentclass[troisieme]{fiche}
\metadonnees{1}{Le terrier}{Bloc A — Raconter}

\begin{document}
\entete

\objectifs{%
\textbf{Langue} : présent simple et présent en \emph{be + -ing} ; troisième personne du
singulier.\par
\textbf{Texte} : Lewis Carroll, \emph{Alice's Adventures in Wonderland} (1865), ch.
\textsc{i}.\par
\textbf{Durée} : 1 h — exercices 20 min, texte et questions 25 min, version 15 min.}

%% ===================================================================
\exercice[12 min]{Présent simple ou présent en \emph{be + -ing}}

\rappel{\textbf{Rappel.} Le présent simple dit ce qui est vrai en général ou ce qui se
répète. Le présent en \emph{be + -ing} dit ce qui se passe maintenant, au moment où l'on
parle.}

\consigne{Mettez le verbe entre parenthèses à la forme qui convient.}

\begin{anglais}
\begin{enumerate}[label=\arabic*.,itemsep=2.5mm,leftmargin=*]
  \item Alice \emph{(sit)} \trou{is sitting} by her sister on the bank.
  \item She \emph{(not like)} \trou{does not like} books without pictures.
  \item Look! A white rabbit \emph{(run)} \trou{is running} across the field.
  \item Rabbits usually \emph{(not wear)} \trou{do not wear} waistcoats.
  \item ``Why \emph{(you / run)} \trou{are you running} so fast?'' asked Alice.
  \item Alice \emph{(feel)} \trou{feels} sleepy every time the day is hot.
  \item The rabbit \emph{(look)} \trou{is looking} at its watch at this very moment.
  \item Children \emph{(play)} \trou{play} in that field every afternoon.
  \item What \emph{(Alice / think)} \trou{does Alice think} of her sister's book?
  \item Listen! Somebody \emph{(come)} \trou{is coming} down the tunnel.
\end{enumerate}
\end{anglais}

\note{Présent simple : 2, 4, 6, 8, 9. Présent en \emph{be + -ing} : 1, 3, 5, 7, 10.}

\note[Les mots qui décident]{\emph{Look!}, \emph{Listen!}, \emph{now}, \emph{at this very
moment} appellent le \emph{be + -ing}. \emph{Usually}, \emph{every afternoon},
\emph{every time} appellent le présent simple. En l'absence de repère, on se demande si
l'action est en train de se faire ou si elle est habituelle.}

%% ===================================================================
\exercice[8 min]{La troisième personne du singulier}

\rappel{\textbf{Rappel.} Au présent simple, le verbe prend un \emph{-s} à la troisième
personne du singulier. À la forme négative et interrogative, c'est \emph{do} ou \emph{does}
qui porte la marque, et le verbe revient à sa base.}

\consigne{Réécrivez la phrase comme demandé.}

\begin{anglais}
\begin{enumerate}[label=\arabic*.,itemsep=3mm,leftmargin=*]
  \item I read every evening. \emph{(she)}
    \lignes{1}
    \corr{She reads every evening.}
  \item We watch the rabbit. \emph{(he)}
    \lignes{1}
    \corr{He watches the rabbit.}
  \item They carry a watch. \emph{(the rabbit)}
    \lignes{1}
    \corr{The rabbit carries a watch.}
  \item You go to school. \emph{(Alice)}
    \lignes{1}
    \corr{Alice goes to school.}
  \item She likes daisies. \emph{(forme négative)}
    \lignes{1}
    \corr{She does not like daisies.}
  \item He runs across the field. \emph{(forme interrogative)}
    \lignes{1}
    \corr{Does he run across the field?}
  \item Alice has a sister. \emph{(forme négative)}
    \lignes{1}
    \corr{Alice does not have a sister.}
  \item The rabbit is late. \emph{(forme interrogative)}
    \lignes{1}
    \corr{Is the rabbit late?}
\end{enumerate}
\end{anglais}

\note[Orthographe du \emph{-s}]{\emph{-es} après \emph{ch}, \emph{sh}, \emph{s}, \emph{x} et
\emph{o} : \emph{watches}, \emph{goes}. Consonne + \emph{y} donne \emph{-ies} :
\emph{carry $\rightarrow$ carries}, mais voyelle + \emph{y} garde le \emph{y} :
\emph{play $\rightarrow$ plays}.}

\note[Le piège de la phrase 5]{Une seule marque de troisième personne par groupe verbal :
\emph{she does not like}, jamais \emph{she does not likes}. Même règle qu'à l'interrogatif :
\emph{does he run?} Phrase 8 : \emph{be} n'a pas besoin de \emph{do}, il s'inverse tout
seul.}

%% ===================================================================
\newpage
\section{Texte}

\noindent\small\textit{Alice s'ennuie au bord d'une rivière, un après-midi d'été, à côté de
sa sœur qui lit.}\normalsize

\begin{texte}
Alice was beginning to get very tired of sitting by her sister on the
\gl[a bank]{bank}{une berge} and of having nothing to do: once or twice she had
\gl[to peep into]{peeped into}{jeter un coup d'œil dans} the book her sister was reading, but
it had no pictures or conversations in it, ``and what is the use of a book,'' thought Alice
``without pictures or conversations?''

So she was considering in her own mind (as well as she could, for the hot day made her feel
very sleepy and stupid), whether the pleasure of making a \gl[a daisy-chain]{daisy-chain}{une
guirlande de pâquerettes} would be worth the trouble of getting up and picking the
\gl[a daisy]{daisies}{une pâquerette}, when suddenly a White Rabbit with pink eyes ran close
by her.

There was nothing so \emph{very} remarkable in that; nor did Alice think it so \emph{very}
much \gl[out of the way]{out of the way}{extraordinaire, hors du commun} to hear the Rabbit
say to itself, ``Oh dear! Oh dear! I shall be late!'' (when she thought it over afterwards, it
occurred to her that she ought to have wondered at this, but at the time it all seemed quite
natural); but when the Rabbit actually \emph{took a watch out of its
\gl[a waistcoat]{waistcoat}{un gilet}-pocket}, and looked at it, and then hurried on, Alice
started to her feet, for it \gl[to flash across]{flashed across}{traverser comme un éclair}
her mind that she had never before seen a rabbit with either a waistcoat-pocket, or a watch to
take out of it, and burning with curiosity, she ran across the field after it, and fortunately
was just in time to see it \gl[to pop down]{pop down}{s'engouffrer dans} a large rabbit-hole
under the \gl[a hedge]{hedge}{une haie}.

In another moment down went Alice after it, never once considering how in the world she was
to get out again.
\end{texte}
\source{Lewis Carroll, \emph{Alice's Adventures in Wonderland}, Londres, Macmillan, 1865,
ch.~\textsc{i}.}

\partiel[15 min]{Questions}
\consigne{Répondez aux deux premières questions en français, aux deux dernières en anglais.}

\begin{enumerate}[label=\arabic*.,itemsep=2.5mm,leftmargin=*]
  \item Pourquoi Alice s'ennuie-t-elle au début du texte ?
    \lignes{2}
    \corr{Elle est assise à ne rien faire à côté de sa sœur, et le livre de celle-ci, dans
    lequel elle a jeté un œil, ne contient ni images ni dialogues.}
  \item Qu'est-ce qui, chez le lapin, décide Alice à le suivre ? Deux choses exactement.
    \lignes{2}
    \corr{Non pas qu'il parle — cela lui a paru naturel sur le moment — mais qu'il sorte une
    montre de la poche de son gilet : elle n'avait jamais vu un lapin avec une poche de
    gilet, ni avec une montre à en sortir.}
  \item \en{What does Alice do as soon as the Rabbit hurries on?}
    \lignes{2}
    \corr{She jumps to her feet, runs across the field after it, and gets there just in time
    to see it pop down a large rabbit-hole under the hedge. Then she goes down after it.}
  \item \en{Find three \emph{-ing} forms in the text. Which two of them follow the verb
    \emph{be}?}
    \lignes{2}
    \corr{\emph{was beginning} et \emph{was reading} (et \emph{was considering}) suivent
    \emph{be}. \emph{Sitting}, \emph{having}, \emph{making}, \emph{getting up},
    \emph{picking}, \emph{burning}, \emph{considering} (le dernier) ne suivent pas
    \emph{be} : ils viennent après une préposition ou décrivent une action simultanée.}
\end{enumerate}

\note[Ce que le texte apprend]{\emph{Was beginning}, \emph{was reading},
\emph{was considering} sont des prétérits en \emph{be + -ing} : la même forme qu'à
l'exercice 1, mais dans le passé. Le prétérit est le sujet de la fiche 2. La phrase
\emph{what is the use of a book?} est au présent parce qu'Alice se le demande à voix haute :
dans un récit au passé, les paroles gardent le temps où elles ont été dites.}

%% ===================================================================
\section{Version}
\consigne{Traduisez le passage en français. Il suit immédiatement le texte.}

\begin{version}
The rabbit-hole went straight on like a tunnel for some way, and then
\gl[to dip]{dipped}{plonger} suddenly down, so suddenly that Alice had not a moment to think
about stopping herself before she found herself falling down a very deep
\gl[a well]{well}{un puits}.

Either the well was very deep, or she fell very slowly, for she had plenty of time as she
went down to look about her and to wonder what was going to happen next.
\end{version}
\source{\emph{Ibid.}}

\lignes{10}

\corr{\textbf{Proposition de traduction.} Le terrier alla tout droit, comme un tunnel, sur
une certaine distance, puis plongea brusquement, si brusquement qu'Alice n'eut pas un instant
pour songer à s'arrêter avant de se retrouver en train de tomber dans un puits très profond.

Ou bien le puits était très profond, ou bien elle tombait très lentement, car elle eut tout
le loisir, en descendant, de regarder autour d'elle et de se demander ce qui allait arriver
ensuite.}

\note[Choix de traduction 1]{\emph{she found herself falling} : \emph{find oneself} +
\emph{-ing} = \og se retrouver en train de \fg. Ne pas traduire \emph{find} par
\og trouver \fg{} tout seul.}

\note[Choix de traduction 2]{\emph{Either\ldots or} : \og ou bien\ldots ou bien \fg. En tête
de phrase, le français répète le \og ou bien \fg, ce que l'anglais ne fait pas.}

\note[Choix de traduction 3]{\emph{what was going to happen next} : \emph{be going to} dans
le passé, donc \og ce qui allait arriver \fg. Le futur proche se rend par \og aller \fg{} +
infinitif, à l'imparfait ici puisque le récit est au passé.}

%% ===================================================================
\partiel{Lexique à mémoriser}
\begin{multicols}{2}\small
\begin{itemize}[label=--,itemsep=0.8mm,leftmargin=*]
  \item \textbf{a bank} — une berge
  \item \textbf{to peep} — jeter un coup d'œil
  \item \textbf{sleepy} — somnolent
  \item \textbf{worth} — qui vaut la peine
  \item \textbf{the trouble} — la peine, le dérangement
  \item \textbf{to pick} — cueillir
  \item \textbf{a hedge} — une haie
  \item \textbf{a waistcoat} — un gilet
  \item \textbf{to hurry} — se dépêcher
  \item \textbf{to wonder} — se demander
  \item \textbf{deep} — profond
  \item \textbf{slowly} — lentement
\end{itemize}
\end{multicols}

\end{document}
